\documentclass[10pt,a4paper]{article}

%% paquetes básicos
\usepackage[utf8]{inputenc}
\usepackage[spanish,es-tabla]{babel}
\usepackage{a4wide,color,verbatim,Sweave,url,xargs,amsmath,booktabs,longtable}
\usepackage{graphicx,float}
\usepackage{tikz,xcolor}
\usepackage{pgfplots}
\usepackage{enumitem}

%% bibliotecas TikZ
\usetikzlibrary{automata,positioning,calc,arrows}
\pgfplotsset{compat=1.18}

%% entornos para exams
\newenvironment{question}{\item}{}
\newenvironment{solution}{\comment}{\endcomment}
\newenvironment{answerlist}{\renewcommand{\labelenumii}{(\alph{enumii})}\begin{enumerate}}{\end{enumerate}}

%% comandos para metadatos exams
\newcommand{\exname}[1]{\def\@exname{#1}}
\newcommand{\extype}[1]{\def\@extype{#1}}
\newcommand{\exsolution}[1]{\def\@exsolution{#1}}
\newcommand{\exshuffle}[1]{\def\@exshuffle{#1}}
\newcommand{\exsection}[1]{\def\@exsection{#1}}

%% configuración párrafos
\setlength{\parskip}{0.7ex plus0.1ex minus0.1ex}
\setlength{\parindent}{0em}

\begin{document}
\Sconcordance{concordance:trigonometria_punto_k_v2.tex:trigonometria_punto_k_v2.Rnw:1 %
35 1 1 146 34 1 1 5 1 1 1 5 1 1 1 5 1 1 1 5 20 1 1 2 4 0 1 2 11 1}


\begin{enumerate}


\begin{question}

Un punto M se mueve de un extremo a otro del segmento PW que se muestra en la gráfica.

% Diagrama geométrico simplificado
\begin{center}
\begin{tikzpicture}[scale=1.0]
% Segmento horizontal QT
\draw[thick] (0,0) -- (4,0);
\node[below] at (0,0) {P};
\node[below] at (4,0) {W};
% Punto K móvil
\fill[blue] (2,0) circle (2pt);
\node[below] at (2,0) {M};
% Punto P fijo
\fill[red] (2,1.5) circle (2pt);
\node[above] at (2,1.5) {Q};
% Línea vertical h
\draw[dashed, red] (2,0) -- (2,1.5);
\node[right] at (2,0.75) {$h$};
% Línea KP
\draw[thick, blue] (2,0) -- (2,1.5);
% Ángulo alpha
\draw[thick] (1.5,0) arc (180:90:0.5);
\node at (1.3,0.3) {$\alpha$};
\end{tikzpicture}
\end{center}

El ángulo $\alpha$ y la medida $h$ se relacionan mediante la razón trigonométrica $\sen(\alpha) = \frac{h}{MQ}$, de donde se deduce la distancia entre M y Q como $MQ = \frac{h}{\sen(\alpha)}$ o $MQ = h \times \csc(\alpha)$.

¿Cuál es la gráfica que muestra las distancias $MQ$, cada vez que M se mueve sobre el segmento PW?

\begin{answerlist}
\item
\includegraphics[width=5cm]{grafica_A.png}
\item
\includegraphics[width=5cm]{grafica_B.png}
\item
\includegraphics[width=5cm]{grafica_C.png}
\item
\includegraphics[width=5cm]{grafica_D.png}
\end{answerlist}

\end{question}

\begin{solution}

La relación matemática fundamental es $MQ = h \times \csc(\alpha) = \frac{h}{\sen(\alpha)}$.

Cuando M se mueve sobre el segmento PW:

\begin{itemize}
\item Al inicio: $\alpha$ es muy pequeño, entonces $\sen(\alpha) \approx 0$, por lo que $\csc(\alpha) \to \infty$ y $MQ$ es muy grande.

\item Al final: $\alpha$ se acerca a 90 grados, entonces $\sen(\alpha) \to 1$, por lo que $\csc(\alpha) \to 1$ y $MQ = h$.
\end{itemize}

Por lo tanto, la gráfica debe mostrar una función cosecante decreciente que empieza con valores muy altos y se aproxima a $h$ cuando $\alpha$ se acerca a 90 grados.

La respuesta correcta es la opción B.

\begin{answerlist}
  \item Falso
  \item Verdadero
  \item Falso
  \item Falso
\end{answerlist}
\end{solution}

%% META-INFORMATION
\exname{Trigonometria Punto Movil}
\extype{schoice}
\exsolution{0100}
\exshuffle{TRUE}
\exsection{Geometria Metrica}

\end{enumerate}
\end{document}
